\documentclass[11pt,a4paper]{article}
\usepackage[T1]{fontenc}
\usepackage[utf8]{inputenc}
\usepackage[dvipsnames]{xcolor}
\usepackage{amsmath,amssymb,amsthm,mathtools}
\usepackage{newpxtext,newpxmath}
\usepackage[a4paper,margin=25mm,headheight=15pt]{geometry}
\usepackage{microtype}
\usepackage{booktabs,longtable,array,multirow,float}
\usepackage{enumitem,fancyhdr,titlesec}
\usepackage{listings}
\usepackage[most]{tcolorbox}
\usepackage[colorlinks=true,linkcolor=MidnightBlue,citecolor=MidnightBlue,urlcolor=MidnightBlue]{hyperref}
\usepackage[nameinlink,noabbrev]{cleveref}
\definecolor{Ink}{HTML}{18324C}
\definecolor{Accent}{HTML}{315D66}
\definecolor{Pale}{HTML}{F0F5F7}
\titleformat{\section}{\Large\bfseries\color{Ink}}{\thesection}{0.65em}{}
\titleformat{\subsection}{\large\bfseries\color{Accent}}{\thesubsection}{0.65em}{}
\titleformat{\subsubsection}{\normalsize\bfseries}{\thesubsubsection}{0.65em}{}
\pagestyle{fancy}
\fancyhf{}
\fancyhead[L]{\small Finite-prefix corrections and dual certificates}
\fancyhead[R]{\small ProveIt research report}
\fancyfoot[C]{\thepage}
\renewcommand{\headrulewidth}{0.3pt}
\setlength{\parindent}{1.3em}
\setlength{\parskip}{0.25em}
\setlist{itemsep=0.25em,topsep=0.4em}
\emergencystretch=2em
\allowdisplaybreaks[2]
\newtheorem{theorem}{Theorem}[section]
\newtheorem{lemma}[theorem]{Lemma}
\newtheorem{proposition}[theorem]{Proposition}
\newtheorem{corollary}[theorem]{Corollary}
\theoremstyle{definition}
\newtheorem{definition}[theorem]{Definition}
\newtheorem{example}[theorem]{Example}
\newtheorem{question}[theorem]{Research question}
\theoremstyle{remark}
\newtheorem{remark}[theorem]{Remark}
\newcommand{\R}{\mathbb R}
\newcommand{\Q}{\mathbb Q}
\newcommand{\N}{\mathbb N}
\newcommand{\Z}{\mathbb Z}
\newcommand{\bA}{\boldsymbol a}
\newcommand{\bX}{\boldsymbol X}
\newcommand{\bone}{\boldsymbol 1}
\newcommand{\spr}{\operatorname{spr}}
\newcommand{\supp}{\operatorname{supp}}
\newcommand{\trunc}[2]{[\xi^{\leq #1}]#2}
\newcommand{\tail}[2]{[\xi^{>#1}]#2}
\newcommand{\repo}{\texttt{ProveIt}}
\newcommand{\sha}{\texttt{e21766d04c2b8a9b2cdba0cd43e563024b3bd1b9}}
\newtcolorbox{resultbox}{colback=Pale,colframe=Accent,boxrule=0.6pt,arc=1mm,left=2.5mm,right=2.5mm,top=2mm,bottom=2mm}
\lstset{basicstyle=\ttfamily\footnotesize,breaklines=true,columns=fullflexible,keepspaces=true,frame=single,rulecolor=\color{gray!40},showstringspaces=false}
\hypersetup{pdftitle={Finite-Prefix Corrections and Dual Certificates for Polyomino Growth},pdfauthor={AI-assisted research report prepared for Vladimir Reshetnikov},pdfsubject={An exact 4.498 upper bound and certificates for convolution recurrences}}

\begin{document}
\hypersetup{pageanchor=false}
\begin{titlepage}
\thispagestyle{empty}
\vspace*{9mm}
\noindent{\color{Accent}\large\scshape ProveIt research series}\par
\vspace{11mm}
\noindent{\color{Ink}\Huge\bfseries Finite-Prefix Corrections\par and Dual Certificates\par for Polyomino Growth}\par
\vspace{7mm}
\noindent{\LARGE An exact bound $\lambda\leq 4.498$\par and a certificate theorem for convolution recurrences}\par
\vspace{10mm}
\noindent{\large Research report prepared for Vladimir Reshetnikov}\par
\noindent September 2026 \quad $\cdot$ \quad AI-assisted research draft\par
\vspace{9mm}
\begin{abstract}
We develop two extensions of the seventeen-neighborhood polyomino project in
\repo. First, exact small-size occurrence counts are incorporated into a
positive convolution system by subtracting its finite coefficient defect.
A nonnegative tail formulation makes this correction rigorous without assuming
convergence of any generating function. Enumerating the seventeen marked
neighborhoods through size $18$ gives an exact rational certificate for
\[
 \boxed{\lambda\leq\frac{2249}{500}=4.498<4.5,}
\]
improving the repository's $4.5235$ upper bound.
Second, balanced integer weights on monomials produce finite arithmetic
certificates of divergence. We prove completeness for productive,
strongly connected, nonlinear positive polynomial systems, thereby providing
a certificate-based answer to Bui's Question~2 for this class. An explicit
$53$-weight certificate and the repository's primal certificate bracket the
uncorrected recurrence growth between $4.52349$ and $4.5235$.
We also prove a spectral stability criterion for the prefix hierarchy and give
a counterexample in which arbitrarily long exact prefixes still fail to
recover the true exponential growth. Full proofs, exact data, enumerators,
rational verifiers, and proposed research directions are supplied.
\end{abstract}
\vfill
\begin{resultbox}
\small
\textbf{Status and scope.} The new bound is a computer-assisted mathematical
result, not a newly kernel-checked Lean theorem. The new arguments are proved
in this report; the finite enumeration and arithmetic are reproducible.
The geometric recurrence lemma is imported from Bui and the pinned repository.
The work has not been independently refereed. General geometric-programming
duality and polyomino enumeration algorithms are prior tools, not claimed
inventions of this report.
\end{resultbox}
\end{titlepage}

\pagenumbering{roman}
\hypersetup{pageanchor=true}
{\small\setlength{\parskip}{0pt}\tableofcontents}
\clearpage
\pagenumbering{arabic}

\section{Problem selection, sources, and main contributions}
\label{sec:scope}

\subsection{Why this repository thread}
The \repo\ repository contains a concrete project on fixed square-lattice
polyominoes at
\begin{center}\small
\texttt{Combinatorics/Polyominoes/KlarnerConstant/}.
\end{center}
Its existing endpoint is an upper bound of $9047/2000=4.5235$ for Klarner's
constant. The relevant sources were read at commit
\begin{center}\small\sha.\end{center}
In particular, the project README, the independent rational verifier,
\texttt{Patterns.lean}, and \texttt{GeometricComplete.lean} fix the mathematical
and computational interface used here \cite{repo,reporeadme,repocert,repopatterns,repocomplete}.
The last of these files assembles the seventeen geometric recurrences for
actual marked polyomino counts and then applies the $4.5235$ certificate.
We inspected those source statements; we did not run a fresh Lean build.

The published input is Vuong Bui's \emph{A convolutional approach to bounding
the number of polyominoes}, arXiv:2511.00461v2, dated 6 May 2026
\cite{bui}. Its Lemma~4 and Appendix~B give the seventeen-neighborhood
recurrence system. Its Section~3, Question~2 asks for compact, readily
verifiable lower-growth certificates for convolution recurrences.
This is a particularly suitable target: the geometric inequalities are
already available, the numerical improvement can have an exact proof
boundary, and the certificate question has a precise published formulation.

\subsection{What is proved here}
The first result concerns actual polyominoes. If $A_n$ is the number of
translation classes of $n$-cell polyominoes, then the data and certificate in
this report give
\begin{equation}\label{eq:headline}
 A_n\leq\frac{927884613}{10^9}\left(\frac{2249}{500}\right)^n
 \quad(n\geq1),\qquad
 \lambda\leq\frac{2249}{500}.
\end{equation}
The factor $927884613/10^9$ is the $G$ coordinate of the supplied rational
supersolution. No numerical eigenvalue or floating-point root appears in the
verification of \eqref{eq:headline}.

The second result concerns the equality recurrences used as majorants.
We give a theorem of alternatives: under explicit productivity,
irreducibility, and nonlinearity hypotheses, either a positive
supersolution exists or a balanced integer-monomial certificate disproves its
existence. The two alternatives cannot occur together. This supplies the
requested type of lower-growth certificate for a class containing Bui's
seventeen-variable system. The explicit example proves
\begin{equation}\label{eq:original-bracket}
 4.52349\leq\mu_{\rm original}\leq4.5235.
\end{equation}
Here $\mu_{\rm original}$ is the exponential growth of the \emph{uncorrected
equality recurrence}, not the polyomino growth constant $\lambda$.
In particular, the lower inequality in \eqref{eq:original-bracket} is
\emph{not} a lower bound on $\lambda$.

The third result describes both the power and a limitation of exact-prefix
corrections. The corrected majorants decrease coefficientwise as more exact
terms are supplied. They eventually admit certificates at any evaluation
point where the linearized system at the true series has spectral radius
less than one. Conversely, a supercritical linearization obstructs the
whole hierarchy under a mild nonzero-tail condition. A scalar example has
true growth $1$ but corrected-majorant growth tending to $3$.

\subsection{Novelty and dependency boundaries}
The proposed research contributions are the explicit $4.498$ polyomino bound,
the finite-prefix application and its exact data, the arithmetic response to
Bui's certificate question, and the accompanying analysis of the prefix
hierarchy. The examined repository snapshot and primary sources did not
contain this $4.498$ certificate or this marked-prefix calculation. That is a
limited novelty assessment, not an exhaustive priority determination.

Several ingredients are established mathematics. The geometric inequalities
come from Bui. Enumeration by canonical include/exclude growth is in the
tradition of Redelmeier \cite{redelmeier}. Positive monomial inequalities,
logarithmic convexification, and entropy-like dual expressions belong to
geometric programming \cite{duffin,boyd,bv}. We do not present those tools as
new discoveries. The proof of the certificate alternative is included to
make its hypotheses, strictness, and rational verification completely
explicit in the recurrence setting.

\begin{resultbox}
\textbf{Three distinct levels.}
The geometric recurrence theorem is an imported mathematical input.
The new prefix and dual-certificate theorems have written proofs here.
The headline numerical bound additionally uses the finite enumeration in
Appendix~\ref{app:data}, whose source programs and exact outputs are supplied.
Arithmetic verification is not a substitute for verifying enumeration.
\end{resultbox}

\section{Polyomino counts and the positive polynomial system}
\label{sec:system}

\subsection{Counting conventions}
A polyomino is a nonempty finite edge-connected subset of $\Z^2$.
Two such subsets are equivalent precisely when one is a translate of the
other. Rotations and reflections are not identified. Holes are allowed.
Let $A_n$ be the number of equivalence classes with $n$ cells, for $n\geq1$.
For the recurrence calculation we put $A_0=0$ and all marked size-zero
counts equal to zero. This convention differs from the optional empty-object
term $A_0=1$ in some sequence tables; positive-size terms are unchanged.

The standard concatenation argument gives $A_{m+n}\geq A_mA_n$ for positive
$m,n$. Fekete's lemma then gives
\begin{equation}\label{eq:lambda}
 \lambda=\lim_{n\to\infty}A_n^{1/n}=\sup_{n\geq1}A_n^{1/n}.
\end{equation}
These are background inputs also implemented in the repository's
\texttt{Concatenation.lean} and \texttt{Asymptotic.lean}
\cite{bui,reporeadme,repocomplete}.

A neighborhood pattern consists of a required offset set $R$ containing
$(0,0)$ and a disjoint forbidden offset set $E$. An occurrence at a cell $a$
of a polyomino $P$ means
\[
 a+R\subseteq P,\qquad (a+E)\cap P=\varnothing.
\]
The marked count for a pattern is the sum of its occurrences over all fixed
polyominoes of the given size. Counting occurrences, rather than merely
polyominoes which have an occurrence, is essential.

Write $C_n,D_n,\ldots,Z_n$ for the seventeen counts, in the coordinate order
\begin{equation}\label{eq:order}
 \mathcal I=(C,D,E,F,G,H,P,Q,R,S,T,U,V,W,X,Y,Z).
\end{equation}
To avoid confusing a coordinate name with a function, the polynomial map
will be denoted by $\Phi$.

\subsection{Exact offset specification}
Let
\[
 B=\{(-1,-1),(0,-1),(1,-1)\},\qquad
 B_4=B\cup\{(2,-1)\}.
\]
In \cref{tab:patterns}, $1,2,3$ in the required column mean respectively
$\{(0,0)\}$, $\{(0,0),(1,0)\}$, and
$\{(0,0),(1,0),(2,0)\}$. A blank lattice position imposes no condition.
These coordinate sets agree with the repository's pattern table and Bui's
page-15 neighborhood diagrams \cite{repopatterns,bui}.

\begin{table}[htbp]
\centering\small
\begin{tabular}{@{}ccl@{}}
\toprule
Type & Required & Forbidden offsets\\
\midrule
$C$ & 1 & $B\cup\{(-1,1),(1,1),(-1,0),(1,0)\}$\\
$D$ & 1 & $B\cup\{(-1,1),(-1,0),(1,0)\}$\\
$E$ & 1 & $B\cup\{(-1,0),(1,0)\}$\\
$F$ & 1 & $B$\\
$G$ & 1 & $B\cup\{(-1,0)\}$\\
$H$ & 1 & $B\cup\{(-1,1),(-1,0)\}$\\
$P$ & 2 & $\{(0,-1),(1,-1),(2,-1)\}$\\
$Q$ & 2 & $B\cup\{(-1,0)\}$\\
$R$ & 2 & $B_4\cup\{(-1,1),(-1,0)\}$\\
$S$ & 2 & $B\cup\{(-1,1),(-1,0)\}$\\
$T$ & 2 & $B_4\cup\{(-1,0)\}$\\
$U$ & 2 & $B_4$\\
$V$ & 3 & $B_4\cup\{(-1,0)\}$\\
$W$ & 3 & $B_4\cup\{(-1,1),(-1,0)\}$\\
$X$ & 2 & $B_4\cup\{(-1,0),(2,0)\}$\\
$Y$ & 2 & $B_4\cup\{(-1,1),(-1,0),(2,0)\}$\\
$Z$ & 2 & $B_4\cup\{(-1,1),(2,1),(-1,0),(2,0)\}$\\
\bottomrule
\end{tabular}
\caption{The seventeen marked neighborhood patterns.}
\label{tab:patterns}
\end{table}

Every polyomino has a $G$ occurrence at the leftmost cell of its bottom row.
There are at most $n$ possible marked cells. Consequently
\begin{equation}\label{eq:gdom}
 A_n\leq G_n\leq nA_n.
\end{equation}
The first six size-one counts equal one; the other eleven equal zero.

\subsection{The imported coefficient inequalities}
Use the formal variable $\xi$ and coordinate vector
$\bX=(c,d,e,f,g,h,p,q,r,s,t,u,v,w,x,y,z)$. Define
\begin{equation}\label{eq:map}
\begin{aligned}
\Phi_C&=\xi+\xi e,&\Phi_D&=\xi+\xi g,&\Phi_E&=\xi+\xi f,\\
\Phi_F&=g+p,&\Phi_G&=e+q,&\Phi_H&=d+s,\\
\Phi_P&=eh+qd+xr+vy+uyz,\\
\Phi_Q&=\xi g+\xi ge+\xi^2u+\xi^2tg+\xi^2ru,\\
\Phi_R&=y+w,\\
\Phi_S&=\xi g+\xi e^2+\xi^2t+\xi^2xg+\xi^2yu,\\
\Phi_T&=x+v,\\
\Phi_U&=dh+sd+yr+wy+uz^2,\\
\Phi_V&=\xi s+\xi^2g^2+\xi^2te+\xi^2rt,\\
\Phi_W&=\xi s+\xi^2eg+\xi^2xe+\xi^2yt,\\
\Phi_X&=\xi d+\xi^2g+\xi^2u,\\
\Phi_Y&=\xi c+\xi^2g+\xi^2t,\\
\Phi_Z&=\xi c+\xi^2e+\xi^2x.
\end{aligned}
\end{equation}
There are $53$ monomials, each with coefficient one. The order displayed
within each row is also the order of the dual weights in
\cref{tab:dual}. In the first two displayed lines, each coordinate has its
own two-term row.

For the true generating series
$\bA_i(\xi)=\sum_{n\geq1}a_{i,n}\xi^n$, Bui's geometric lemma is the
coefficientwise inequality
\begin{equation}\label{eq:imported}
 \bA\preceq\Phi(\xi,\bA).
\end{equation}
All products in this equation are formal Cauchy products. Thus convergence is
not part of the statement. Equations with a factor $\xi^2$, for example,
encode a two-cell shift in the corresponding convolution index.

The geometric proof of \eqref{eq:imported} uses occupancy partitions,
injective leaf deletions, and decompositions into connected marked pieces.
We import that lemma rather than silently treating a numerical recurrence
fit as a geometric proof. For example, the $C$ row deletes its marked cell,
whose only possible neighbor is above it. The $P$ and $U$ rows encode
multi-piece decompositions; these are not inferred from the finite table.
The repository supplies their separate finite-injection modules
\cite{bui,repocomplete}.

The five same-size partitions are in fact exact:
\begin{equation}\label{eq:partitions}
 F=G+P,\quad G=E+Q,\quad H=D+S,\quad R=Y+W,\quad T=X+V.
\end{equation}
For example, a type-$G$ anchor either has no right neighbor, yielding type
$E$, or has a right neighbor, yielding type $Q$. Similar left/right occupancy
splits give the other identities. They will serve as independent consistency
checks on the occurrence table.

\section{Exact-prefix correction: the general theorem}
\label{sec:prefix}

\subsection{Proper positive polynomial systems}
\begin{definition}
A \emph{proper positive polynomial system} is a map
$\Phi\in\R_{\geq0}[\xi,X_1,\ldots,X_d]^d$ such that
$\Phi(0,0)=0$ and the nonnegative matrix
\[
 L=\Phi_{\bX}(0,0)
\]
is nilpotent. Here and below vector inequalities are coordinatewise.
\end{definition}

Nilpotence handles same-size linear dependencies. In \eqref{eq:map}, the
only such dependencies are the five rows in \eqref{eq:partitions}; their
directed graph has no cycle. All remaining occurrences of an unknown are
multiplied either by a positive power of $\xi$ or by another unknown.
Thus the system is proper. One possible degree-by-degree evaluation order is
\[
 C,D,E,Q,S,V,W,X,Y,Z,P,U,G,F,H,T,R.
\]
The formal development does not require this particular order.

\begin{lemma}[Formal fixed points and comparison]\label{lem:formal}
A proper positive polynomial system has a unique fixed point in
$\xi\R[[\xi]]^d$. Its coefficients are nonnegative, and the iterates from
zero increase coefficientwise and stabilize at each finite degree.
If $U\in\xi\R_{\geq0}[[\xi]]^d$ satisfies
$U\preceq\Phi(\xi,U)$, then $U$ is coefficientwise at most this fixed point.
\end{lemma}
\begin{proof}
Write $\Phi=L\bX+H$. Because the unknown series have zero constant terms,
the coefficient of degree $n$ in $H(\xi,\bX)$ depends only on unknown
coefficients of degrees strictly less than $n$. If these have already been
chosen, the coefficient vector $b_n$ in a fixed point must satisfy
\[
 (I-L)b_n=h_n.
\]
For some $k$, $L^k=0$, so
\[
 (I-L)^{-1}=I+L+\cdots+L^{k-1}\geq0.
\]
This proves existence, uniqueness, and nonnegativity by induction on $n$.
For a subsolution, the same argument gives
$(I-L)u_n\leq h_n(u_1,\ldots,u_{n-1})$; multiplication by the nonnegative
inverse and induction prove comparison.

The iterates from zero are increasing by monotonicity. If all lower-degree
coefficients have stabilized, the remaining degree-$n$ error is multiplied
by $L$ at each further iteration, and therefore vanishes after at most $k$
steps. This proves finite-degree stabilization.
\end{proof}

The subsolution conclusion uses formal properness. It would be false for
arbitrary numerical subsolutions: $2\leq2^2$ does not imply that $2$ is below
the least nonnegative fixed point of $x\mapsto x^2$.

\subsection{The finite defect and the nonnegative tail map}
Suppose $\bA\in\xi\R_{\geq0}[[\xi]]^d$ is a true counting series satisfying
\eqref{eq:imported}. For $N\geq1$, set
\begin{equation}\label{eq:prefix-def}
 P_N=\trunc{N}{\bA},\qquad
 T_N=\trunc{N}{\Phi(\xi,P_N)},\qquad
 D_N=T_N-P_N.
\end{equation}
The coefficient of $\Phi(\xi,\bA)$ at a degree at most $N$ depends only on
$P_N$, so $D_N\succeq0$. It is precisely the finite prefix of the true
recurrence defect $\Phi(\xi,\bA)-\bA$.

The signed equation
\begin{equation}\label{eq:signed}
 B_N=\Phi(\xi,B_N)-D_N
\end{equation}
is a convenient shorthand, but its right-hand side is not a positive
polynomial in $\xi$. The correct positivity argument uses $Y_N=B_N-P_N$
and the map
\begin{equation}\label{eq:psi}
\begin{split}
 \Psi_N(\xi,Y)
 &=\Phi(\xi,P_N+Y)-T_N\\
 &=\Phi(\xi,P_N+Y)-\Phi(\xi,P_N)
   +\tail{N}{\Phi(\xi,P_N)}.
\end{split}
\end{equation}
Every coefficient of $\Psi_N$ in $\xi,Y_1,\ldots,Y_d$ is nonnegative.
Indeed, in the difference on the second line, the terms independent of $Y$
cancel, leaving only monomials containing at least one $Y$.
Moreover, $\Psi_N(0,0)=0$ and $(\Psi_N)_Y(0,0)=L$.

\begin{theorem}[Prefix-corrected majorant]\label{thm:majorant}
The map $\Psi_N$ has a unique fixed point
$Y_N\in\xi\R_{\geq0}[[\xi]]^d$. Its coefficients through degree $N$ vanish.
The series $B_N=P_N+Y_N$ satisfies
\[
 \bA\preceq B_N,\qquad
 \trunc{N}{B_N}=P_N,
\]
and solves \eqref{eq:signed}.
\end{theorem}
\begin{proof}
The fixed-point statement follows from \cref{lem:formal}. The constant-in-$Y$
part of $\Psi_N$ begins above degree $N$, so degree induction gives
$[\xi^n]Y_N=0$ for $n\leq N$.
Put $U=\bA-P_N\succeq0$. The low-degree part of
$\Phi(\xi,\bA)-\bA$ is $D_N$, and its remaining coefficients are
nonnegative. Therefore
\[
 U\preceq\Phi(\xi,P_N+U)-T_N=\Psi_N(\xi,U).
\]
Formal comparison yields $U\preceq Y_N$. Adding $P_N$ proves domination.
Substituting the fixed-point equation for $Y_N$ yields \eqref{eq:signed}.
\end{proof}

\begin{theorem}[Rational upper certificate]\label{thm:upper}
Let $\zeta>0$. Suppose a vector $v\in\R^d$ satisfies
\begin{equation}\label{eq:upper-conditions}
 v\geq P_N(\zeta),\qquad
 v\geq\Phi(\zeta,v)-D_N(\zeta).
\end{equation}
Then all coordinates of $B_N$ converge at $\zeta$, and
\[
 B_N(\zeta)\leq v,\qquad
 a_{i,n}\leq v_i\zeta^{-n}.
\]
For the polyomino system this implies
$A_n\leq v_G\zeta^{-n}$ and $\lambda\leq\zeta^{-1}$.
When the input prefix, $\zeta$, and $v$ are rational, verification of
\eqref{eq:upper-conditions} requires rational arithmetic only.
\end{theorem}
\begin{proof}
Let $w=v-P_N(\zeta)\geq0$. Equations \eqref{eq:prefix-def} and
\eqref{eq:psi} show that \eqref{eq:upper-conditions} is equivalent to
\[
 \Psi_N(\zeta,w)\leq w.
\]
The numerical iterates $w^{(0)}=0$,
$w^{(j+1)}=\Psi_N(\zeta,w^{(j)})$ therefore stay between $0$ and $w$.
The corresponding formal iterates are polynomials with nonnegative
coefficients. Their evaluation at $\zeta$ is exactly $w^{(j)}$.
For every fixed degree $M$, a sufficiently late formal iterate agrees with
$Y_N$ through degree $M$ by \cref{lem:formal}. Hence
\[
 \sum_{n=0}^M[\xi^n]Y_N\,\zeta^n\leq w.
\]
Monotone convergence of these bounded finite sums proves
$Y_N(\zeta)\leq w$. Adding $P_N(\zeta)$ proves the asserted bound for $B_N$.
Every nonnegative term of a bounded series is bounded by its sum, yielding
the coefficient estimate. Now use \eqref{eq:gdom} and \eqref{eq:lambda}.
\end{proof}

\begin{remark}[The tail condition is essential]
A solution of $v\geq\Phi(\zeta,v)-D_N(\zeta)$ alone need not be a valid
certificate. The subtraction can create extraneous branches or vectors below
the known prefix. The separate inequality $v\geq P_N(\zeta)$ is part of the
proof and is checked explicitly by the verifier.
\end{remark}

\subsection{Monotonicity of the hierarchy}
\begin{theorem}[More exact data never worsens the majorant]\label{thm:nested}
For every $N\geq1$,
\[
 \bA\preceq B_{N+1}\preceq B_N.
\]
Thus the radii of convergence of the corrected majorants are nondecreasing,
and their reciprocal exponential-growth bounds are nonincreasing.
Any real upper certificate for prefix $N$ is also an upper certificate for
prefix $N+1$ at the same evaluation point.
\end{theorem}
\begin{proof}
The first inequality is \cref{thm:majorant}. Through degree $N$, the two
majorants agree. At degree $N+1$, $B_{N+1}$ has the true coefficient and
$B_N$ has a coefficient at least as large. Above that degree, both obey the
same positive coefficient recurrence, because the defect polynomials have
no terms there. Degree induction using $(I-L)^{-1}\geq0$ proves the second
inequality. Coefficientwise domination implies the radius assertion.

If $v$ is an $N$-certificate, then the preceding theorem and
$\bA\preceq B_N$ show that $P_{N+1}(\zeta)\leq v$.
Also $D_{N+1}-D_N$ has nonnegative coefficients, so
$\Phi(\zeta,v)-D_{N+1}(\zeta)\leq v$.
Both certificate conditions for $N+1$ follow.
\end{proof}

For the present system the defect is zero through size $6$. Its size-$7$
coefficient, in the order \eqref{eq:order}, is
\begin{equation}\label{eq:first-defect}
 [\xi^7]D_7=(0,1,2,0,0,0,0,0,0,1,0,0,0,0,0,1,2).
\end{equation}
For example, $a_{D,7}=260$, whereas the deletion bound uses
$a_{G,6}=261$; the corresponding defect coefficient is $1$.
This first discrepancy explains why prefix lengths $1$ through $6$ do not
improve the original equality system.

\section{An exact upper bound below 4.5}
\label{sec:numeric}

\subsection{The certificate}
Take $N=18$ and
\[
 \zeta=\frac{500}{2249}.
\]
Use the exact marked counts in Appendix~\ref{app:data} to construct
$P_{18},T_{18},D_{18}$ by \eqref{eq:prefix-def}.
The vector $v$ is given in \cref{tab:upper}; every numerator in its second
column has common denominator $10^9$.

Define the residual and tail budget by
\begin{equation}\label{eq:slacks}
 \sigma_i=v_i-\Phi_i(\zeta,v)+D_{18,i}(\zeta),\qquad
 \eta_i=v_i-P_{18,i}(\zeta).
\end{equation}
The last two columns are conservative lower bounds:
\[
 \sigma_i\geq\frac{\text{third-column entry}}{10^{12}},\qquad
 \eta_i\geq\frac{\text{fourth-column entry}}{10^6}.
\]
They are obtained by flooring the exact rational numbers after scaling;
no printed decimal approximation is being used to decide a sign.

\begin{table}[htbp]
\centering\small
\begin{tabular}{@{}lrrr@{}}
\toprule
Coordinate & Numerator of $v_i$ & $10^{12}$-scaled slack & $10^6$-scaled tail\\
\midrule
$c$ & 347,752,792 & 99,824 & 28,076 \\
$d$ & 427,615,666 & 100,303 & 59,180 \\
$e$ & 564,832,205 & 100,345 & 123,441 \\
$f$ & 1,551,872,652 & 100,000 & 542,918 \\
$g$ & 927,884,613 & 99,000 & 260,176 \\
$h$ & 727,305,121 & 100,000 & 159,091 \\
$p$ & 623,987,939 & 99,749 & 282,741 \\
$q$ & 363,052,309 & 100,459 & 136,735 \\
$r$ & 234,092,913 & 100,000 & 63,133 \\
$s$ & 299,689,355 & 99,657 & 99,910 \\
$t$ & 283,350,363 & 100,000 & 89,179 \\
$u$ & 487,086,219 & 100,056 & 192,950 \\
$v$ & 119,025,613 & 99,875 & 52,209 \\
$w$ & 97,883,000 & 99,767 & 38,611 \\
$x$ & 164,324,650 & 100,195 & 36,970 \\
$y$ & 136,209,813 & 99,861 & 24,521 \\
$z$ & 112,439,657 & 100,134 & 14,724 \\
\bottomrule
\end{tabular}

\caption{The rational size-$18$ upper certificate and exact lower bounds on
its margins. All vector entries have denominator $10^9$.}
\label{tab:upper}
\end{table}

\begin{theorem}[Computer-assisted polyomino bound]\label{thm:main}
For the fixed-polyomino counting convention in \cref{sec:system},
\[
 A_n\leq\frac{927884613}{10^9}
          \left(\frac{2249}{500}\right)^n\quad(n\geq1),
 \qquad \lambda\leq4.498<4.5.
\]
\end{theorem}
\begin{proof}
The finite enumeration described below produces the true coefficients of
$P_{18}$ in Appendix~\ref{app:data}. Bui's recurrence lemma gives
\eqref{eq:imported}. Exact integer polynomial convolution constructs $D_{18}$.
Substituting the rational vector of \cref{tab:upper} into
\eqref{eq:slacks} gives
\[
 \min_i\sigma_i\geq\frac{99}{10^9}>0,\qquad
 \min_i\eta_i\geq\frac{14724}{10^6}>0.
\]
Thus \cref{thm:upper} applies. Its $G$ coordinate is
$v_G=927884613/10^9$, yielding the claimed inequalities.
\end{proof}

The strict residuals also show that the certificate is not balanced on a
rounding boundary. By continuity, the same vector works for a sufficiently
small increase of $\zeta$. We retain $2249/500$ as the stated bound because
it is simple and explicitly verified.

\subsection{Discovery is not verification}
The search for a candidate used floating-point solutions of
\[
 \Phi(\zeta,v)-D_N(\zeta)=v,
 \qquad \det(I-\Phi_{\bX}(\zeta,v))=0.
\]
These equations suggest where the equality envelope becomes critical.
For a slightly smaller evaluation point, we solved a perturbed fixed-point
equation with a small positive residual in every coordinate, then rounded the
candidate to denominator $10^9$. Such a solver may converge to an undesired
branch or misestimate a threshold. Neither event can invalidate an accepted
certificate: the verifier checks the two inequalities in
\eqref{eq:upper-conditions} independently from scratch.

\Cref{tab:progress} records the observed progression. The middle column
contains numerical critical-point estimates, not proved optimal values or
certified lower bounds. Every entry in the last column has a separate
exact rational certificate in the archive.

\begin{table}[htbp]
\centering\small
\begin{tabular}{@{}rrr@{}}
\toprule
Prefix size $N$ & Estimated equality growth & Certified upper bound\\
\midrule
$1$--$6$ & $4.523499228$ & $4.5235$ (original certificate)\\
7 & $4.523230947$ & $4.5233$\\
8 & $4.522434056$ & $4.5225$\\
9 & $4.521121227$ & $4.5212$\\
10 & $4.519368978$ & $4.5194$\\
11 & $4.517262575$ & $4.5173$\\
12 & $4.514881253$ & $4.5149$\\
13 & $4.512294495$ & $4.5123$\\
14 & $4.509560290$ & $4.5096$\\
15 & $4.506725656$ & $4.5068$\\
16 & $4.503828169$ & $4.5039$\\
17 & $4.500897596$ & $4.5009$\\
18 & $4.497957351$ & $4.4980$\\
\bottomrule
\end{tabular}

\caption{Progress under exact-prefix correction. Numerical estimates and
rigorous upper certificates are deliberately separated.}
\label{tab:progress}
\end{table}

The original recurrence's growth is already within $10^{-5}$ of the
repository's bound, as \cref{sec:dual-example} will prove. Consequently,
\cref{thm:main} does not arise from finding a better supersolution for the
unchanged system. The correction changes the equality envelope while
preserving its validity as an upper bound on the true counts.

\subsection{A reproducible finite enumeration}
Every translation class has a unique representative whose leftmost cell on
the bottom row is $(0,0)$. All its cells then lie in
\[
 \mathcal H=\{(x,y):y>0\}\cup\{(x,0):x\geq0\}.
\]
The main enumerator maintains a connected occupied set $S$, a finite ordered
frontier $U$ of available adjacent cells, and an excluded set $E$.
Initially $S=\varnothing$, $U=\{(0,0)\}$, and $E=\varnothing$.
At a call, take the last frontier cell $p$, include it in $S$, record the
resulting object, and recurse after adding newly available neighbors.
After that branch returns, remove $p$ and exclude it for all later branches
of the same call. The exclusions introduced by a call are removed when that
call returns. No quotient by rotations or reflections is performed.

\begin{proposition}[Enumeration correctness]\label{prop:enumeration}
The include/exclude frontier procedure records each nonempty finite
edge-connected subset of $\mathcal H$ containing $(0,0)$ exactly once, up to
the specified size limit. Thus its totals are precisely the fixed-polyomino
counts.
\end{proposition}
\begin{proof}
Every recorded set is connected because each added cell is adjacent to the
previous occupied set, except for the initial origin. It stays in
$\mathcal H$ by the admissibility test.

Fix a target set $P$ of the required kind. At any intermediate state with
$S\subsetneq P$ and no excluded cell in $P$, connectedness implies that some
cell of $P\setminus S$ lies in the available frontier. In the order processed
by the current call, there is a unique first frontier cell belonging to $P$.
All cells processed before it can be excluded without affecting $P$.
The branch including that unique cell preserves the same invariant for $P$.
Induction reaches $S=P$, which is recorded.

The same argument proves uniqueness: the first processed frontier cell
belonging to $P$ determines the only branch that can contain $P$ without
excluding one of its cells. Applying this observation recursively determines
a unique path. Finally, the chosen anchor gives exactly one representative
of each translation class.
\end{proof}

For an occupied anchor, its membership in each pattern depends only on ten
nearby positions:
\[
 (-1,1),(1,1),(2,1),(-1,0),(1,0),(2,0),
 (-1,-1),(0,-1),(1,-1),(2,-1).
\]
These are encoded by a ten-bit mask. All $2^{10}$ masks are preclassified
against the required and forbidden sets of \cref{tab:patterns}.
The direct implementation scans all occupied anchors whenever an object is
recorded. The faster implementation maintains the pattern totals under cell
insertion and deletion. Inserting $p$ changes only the contribution of $p$
and those anchors $p-o$ for which $o$ is one of the ten offsets. Removing $p$
reverses exactly these changes. This local invariant proves that incremental
updates give the same occurrence count as a complete scan.

The size-$18$ computation records
\[
 A_{18}=1,540,820,542,\qquad
 \sum_{n=1}^{18}A_n=2,083,404,030
\]
objects across all recorded sizes. The unmarked totals through $18$ agree
with OEIS A001168 \cite{oeis}; those familiar totals are a consistency check,
not a replacement for the marked enumeration.

\subsection{Finite-machine safety and independent checks}
The C++ sources are restricted to size at most $18$. Any cell of an anchored
$n$-cell polyomino satisfies $|x|\leq n-1$ and $0\leq y\leq n-1$.
The padded array is wide and tall enough for these coordinates and for all
mask and neighbor offsets. Thus no boundary wrapping or out-of-range
access is needed by the algorithm.

The counters are unsigned $64$-bit integers. An a priori bound verifies that
overflow is impossible at these sizes. A deterministic depth-first spanning
tree of an anchored $n$-cell set can be encoded by a plane rooted tree and
directions for its $n-1$ forward steps. There are
$\operatorname{Cat}_{n-1}$ plane tree shapes. The first forward step has at
most four choices; every later forward step has at most three, since a
nonroot vertex cannot move back to its parent and the root has already used
one child direction. Consequently, for $n\geq2$,
\[
 A_n\leq4\cdot3^{n-2}\operatorname{Cat}_{n-1}.
\]
This overcounts heavily but is sufficient. Every marked count is at most
$nA_n$, and the largest bound for $n\leq18$ is
\[
 18\cdot4\cdot3^{16}\operatorname{Cat}_{17}
 =401816383504818480<2^{64}.
\]
Incremental counters count occurrences in the currently occupied finite set;
the insertion/deletion invariant also prevents underflow.

Two structurally different enumeration approaches were compared through size
$11$: the C++ include/exclude procedure and a Python procedure that repeatedly
adds each boundary cell, translates each shape to a canonical coordinate
tuple, and deduplicates using a set. All $187$ marked values and all unmarked
totals agreed. Direct scanning and incremental scanning in C++ agreed through
size $16$. The final size-$18$ output also agrees with the previous
size-$16$ table on every earlier entry.

These checks do not amount to independent formal certification of all
size-$18$ values. They locate the remaining computational trust boundary:
the correctness of the supplied enumerator, its implementation, and its
execution. The recurrence inequalities themselves are not learned from this
table. Five exact partition identities and all nonnegative recurrence defects
are checked again by the rational verifier.

\section{Balanced-monomial certificates of divergence}
\label{sec:dual}

\subsection{The certificate format}
Consider a finite positive polynomial system
\begin{equation}\label{eq:general-map}
 \Phi_i(\xi,X)=\sum_{k\in K_i}c_{ik}\xi^{b_{ik}}X^{\alpha_{ik}},
 \quad c_{ik}\in\Q_{>0},\quad b_{ik}\in\N,
 \quad \alpha_{ik}\in\N^d.
\end{equation}
Fix a rational $\zeta>0$. We wish to certify that there is no
$X\in\R_{>0}^d$ satisfying $\Phi(\zeta,X)\leq X$.

Assign a nonnegative integer weight $m_{ik}$ to each monomial, not all zero,
and define row sums
\[
 r_i=\sum_{k\in K_i}m_{ik}.
\]
The weights are \emph{balanced} if
\begin{equation}\label{eq:balance}
 \sum_{i,k}m_{ik}\alpha_{ik,j}=r_j\qquad(1\leq j\leq d).
\end{equation}
Define the positive rational product
\begin{equation}\label{eq:dual-product}
 \mathcal C(\zeta,m)=
 \prod_{i,k:m_{ik}>0}
 \left(\frac{c_{ik}\zeta^{b_{ik}}r_i}{m_{ik}}\right)^{m_{ik}}.
\end{equation}
A row of weight zero contributes no factor. The entire certificate consists
of $\zeta$, the integer weights, and a verification that
$\mathcal C(\zeta,m)>1$.

\begin{theorem}[Balanced-monomial obstruction]\label{thm:dual}
If \eqref{eq:balance} holds and $\mathcal C(\zeta,m)>1$, then no positive
supersolution of $\Phi(\zeta,X)\leq X$ exists.
\end{theorem}
\begin{proof}
Suppose such a supersolution exists. For each row with $r_i>0$, put
$p_{ik}=m_{ik}/r_i$ on the positive-weight terms. Weighted arithmetic--geometric
mean gives
\begin{align*}
 1\geq\frac{\Phi_i(\zeta,X)}{X_i}
 &\geq\sum_{k:m_{ik}>0}
        c_{ik}\zeta^{b_{ik}}X^{\alpha_{ik}-e_i}\\
 &\geq\prod_{k:m_{ik}>0}
 \left(\frac{c_{ik}\zeta^{b_{ik}}
                  X^{\alpha_{ik}-e_i}}{p_{ik}}\right)^{p_{ik}}.
\end{align*}
Raise the inequality to $r_i$ and multiply over rows. The exponent of $X_j$
in the resulting right-hand side is
$\sum_{i,k}m_{ik}\alpha_{ik,j}-r_j=0$ by balance.
All unknown variables cancel, leaving
$1\geq\mathcal C(\zeta,m)>1$, a contradiction.
\end{proof}

The cancellation is the essential point. A certificate does not need to
supply a very large iterate, an approximate divergent orbit, or an
approximate critical point. Its inequalities quantify over every positive
candidate supersolution at once.

\begin{example}[Catalan recurrence]
For $\Phi(\xi,x)=\xi+x^2$, assign weight one to both terms.
Then the incoming exponent is $0+2=2=r$, and
\[
 \mathcal C(\zeta,m)=(2\zeta)(2)=4\zeta.
\]
Thus $\zeta>1/4$ immediately has a two-weight divergence certificate, recovering
the usual Catalan threshold from one arithmetic--geometric mean inequality.
\end{example}

\subsection{From nonexistence to a lower growth bound}
Let $B$ be the least formal fixed point of a proper positive system. If its
coordinates converge at $\zeta$, finite products of nonnegative convergent
series can be evaluated, giving
\[
 B(\zeta)=\Phi(\zeta,B(\zeta)).
\]
If the system is productive in the sense defined in the next section, this
fixed point is strictly positive. Therefore \cref{thm:dual} implies that at
least one coordinate fails to converge at $\zeta$.

Define the system's exponential growth by
\[
 \mu=\max_i\limsup_{n\to\infty}\bigl([\xi^n]B_i\bigr)^{1/n}.
\]
Equivalently, $1/\mu$ is the minimum of the coordinate radii of convergence.
Cauchy--Hadamard then gives
\begin{equation}\label{eq:dual-growth}
 \mu\geq\zeta^{-1}.
\end{equation}
We use the limsup formulation so that no unproved assertion about periodicity
or existence of individual root limits is needed.

For productive systems with a strongly connected dependency graph, the
coordinate radii coincide. Indeed, for every dependency edge $i\to j$,
choose one positive coefficient in each of the other factors of a monomial
containing $X_j$. The formal fixed-point equation implies
$B_i\succeq c\xi^kB_j$ for some $c>0$ and integer $k\geq0$.
Thus $R_i\leq R_j$. Following paths in both directions proves equality.

\begin{resultbox}
\textbf{Do not reverse the majorant inequality.}
The true polyomino series is below the equality solution, so an upper bound
on equality growth bounds $\lambda$ from above. A lower bound on equality
growth does not bound $\lambda$ from below. The dual certificate is useful
for certifying the limitation of a chosen recurrence model and for answering
the recurrence-growth question, not for reversing this implication.
\end{resultbox}

\section{Completeness of the certificate alternative}
\label{sec:complete}

\subsection{The hypotheses}
For this section fix $\zeta>0$ and write $F(X)=\Phi(\zeta,X)$.
The coefficients of $F$ are positive on its listed monomials.

\begin{definition}
The map $F$ is \emph{productive} if some finite iterate of $F$, starting from
zero, has every coordinate strictly positive. Its dependency graph has an
edge $i\to j$ if some monomial of $F_i$ contains $X_j$.
The map is \emph{strongly connected} if this graph is strongly connected,
and \emph{nonlinear} if at least one monomial has total degree at least two
in the unknowns.
\end{definition}

Productivity is stronger than merely having a constant term somewhere.
For example, $F_1=1+X_1X_2$, $F_2=X_1X_2$ has a strongly connected graph,
but iteration from zero never makes $X_2$ positive. This degeneracy is
excluded. All three conditions are finite structural checks on the monomial
list; they hold for \eqref{eq:map} and are replayed by the verifier.

\subsection{Compactness of bounded multiplicative residuals}
\begin{lemma}\label{lem:compact}
If $F$ is productive, strongly connected, and nonlinear, then, for every
$K>0$, the set
\[
 S_K=\{X\in\R_{>0}^d:F(X)\leq KX\}
\]
is contained in a compact box $[a,b]\subset\R_{>0}^d$ and is closed in that
box. In particular, it is compact, possibly empty.
\end{lemma}
\begin{proof}
Every $X\in S_K$ satisfies $X\geq K^{-1}F(X)$. Iteration from zero of the
map $K^{-1}F$ therefore stays below $X$. Multiplication by a positive scalar
does not change which coordinates become positive at any iteration, so
productivity supplies a fixed vector $a>0$ with $X\geq a$ for every
$X\in S_K$.

For each edge $i\to j$, choose a monomial containing $X_j$. Lower-bound all
its other factors, including extra copies of $X_j$ if necessary, by the
corresponding coordinates of $a$. This gives a constant $c_{ij}>0$ such that
\[
 F_i(X)\geq c_{ij}X_j\quad(X\geq a).
\]
On $S_K$ it follows that $X_i\geq(c_{ij}/K)X_j$.
Strong connectivity now gives constants $u_i,v_i>0$ with
\[
 u_iX_1\leq X_i\leq v_iX_1
 \qquad(X\in S_K).
\]
Choose a nonlinear monomial of total degree $q\geq2$ in row $h$.
The lower comparisons give $F_h(X)\geq cX_1^q$ for some $c>0$,
whereas $F_h(X)\leq KX_h\leq Kv_hX_1$ on $S_K$.
Thus $X_1^{q-1}\leq Kv_h/c$, a uniform upper bound.
The other coordinates are then bounded above by their comparisons with
$X_1$. Finally, the polynomial inequalities are closed within this positive
compact box.
\end{proof}

\subsection{An elementary convex proof}
Put $X_j=e^{y_j}$ and define
\begin{equation}\label{eq:h}
 f_i(y)=\log F_i(e^y)-y_i,
 \qquad h(y)=\max_i f_i(y).
\end{equation}
Each $f_i$ is a log-sum-exp of affine functions and is convex. For
completeness, its Hessian is the covariance matrix of the monomial exponent
vectors under their positive softmax probabilities, and is therefore
positive semidefinite. Sublevel sets of $h$ correspond to the sets $S_K$ of
\cref{lem:compact}. Hence $h$ attains its global minimum.

\begin{theorem}[Complete rational certificate alternative]\label{thm:alternative}
Suppose $F$ has rational coefficients and is productive, strongly connected,
and nonlinear. Exactly one of the following holds:
\begin{enumerate}[label=(\roman*)]
 \item there is $X>0$ with $F(X)\leq X$;
 \item there are balanced nonnegative integer monomial weights, not all zero,
       for which the product \eqref{eq:dual-product} is greater than one.
\end{enumerate}
Here the coefficient $c_{ik}\zeta^{b_{ik}}$ in
\eqref{eq:dual-product} is simply the corresponding coefficient of $F$.
\end{theorem}
\begin{proof}
Mutual exclusion is \cref{thm:dual}. Suppose (i) fails. Let $y^*$ minimize
$h$, and put $\delta=h(y^*)$. Since a minimum is attained and no point has
$h\leq0$, we have $\delta>0$.

Let $I$ be the active rows, those with $f_i(y^*)=\delta$.
The origin belongs to the convex hull of their gradients. To see this
without invoking an optimization duality theorem, suppose otherwise and let
$g$ be the point of that compact convex hull nearest the origin.
Then $g\ne0$ and
$g\cdot\nabla f_i(y^*)\geq\|g\|^2$ for every active row.
A sufficiently small step in direction $-g$ strictly decreases every active
$f_i$. The inactive rows remain below the old maximum by continuity and
finiteness. This contradicts minimality of $h$.
Consequently there are numbers $\theta_i\geq0$, supported on $I$, such that
\[
 \sum_i\theta_i=1,\qquad
 \sum_i\theta_i\nabla f_i(y^*)=0.
\]

For each monomial let
\[
 p_{ik}=
 \frac{c_{ik}\zeta^{b_{ik}}e^{\alpha_{ik}\cdot y^*}}
      {F_i(e^{y^*})},\qquad
 \beta_{ik}=\theta_i p_{ik}.
\]
Then $\sum_k\beta_{ik}=\theta_i$, and the gradient identity is exactly
\[
 \sum_{i,k}\beta_{ik}\alpha_{ik,j}=\theta_j.
\]
Thus the real weights $\beta$ are balanced and have total mass one.
Their logarithmic product is
\begin{align*}
 \mathcal L(\beta)
 &=\sum_{i,k:\beta_{ik}>0}
 \beta_{ik}\log\left(
   \frac{c_{ik}\zeta^{b_{ik}}\theta_i}{\beta_{ik}}\right)\\
 &=\sum_i\theta_i\log F_i(e^{y^*})
   -\sum_j\theta_j y_j^*\\
 &=\sum_i\theta_i f_i(y^*)=\delta>0.
\end{align*}

It remains to replace real weights by integers. Restrict to the support of
$\beta$. On this support, balance and total mass one form a consistent
linear system with rational coefficients. Gaussian elimination supplies a
rational particular solution and a rational basis for the homogeneous
solution space. Approximating the real free parameters by rationals gives
balanced rational weights arbitrarily close to $\beta$, positive on the
same support. The displayed logarithmic expression is continuous there,
so a sufficiently close rational approximation still has positive value.
Multiplying by a common denominator gives integer weights. Scaling multiplies
the logarithmic product by that positive denominator, preserving strict
positivity. This proves (ii).
\end{proof}

\subsection{What this settles, and what it does not}
A positive numerical supersolution bounds the zero iterates and hence all
formal partial sums, by the same argument as in \cref{thm:upper}. Thus it
forces convergence at the evaluation point. For a proper productive system
in the stated class, divergence at a rational point therefore has a finite
balanced-monomial obstruction, and a finite obstruction
implies the lower-growth bound \eqref{eq:dual-growth}. Thus there is a
certificate alternative, not just a test based on watching a floating-point
iteration grow. This answers the certificate-existence aspect of Bui's
Question~2 for this class, including his seventeen-variable system.

The result does not assert that the original question for every conceivable
convolution, maximum, or signed recurrence has been settled. Nor does it give
a uniform polynomial bound on the bit length of the weights as a target
approaches the exact growth threshold. The explicit certificate below is
small and quickly checked; a worst-case certificate-size theory remains a
separate problem. The underlying convex mechanism is classical geometric
programming, here specialized and proved in an arithmetic certificate form
\cite{duffin,boyd,bv}.

\section{A compact dual certificate for the original recurrence}
\label{sec:dual-example}

\subsection{The integer weights}
For the uncorrected map \eqref{eq:map}, take
\[
 \zeta_{\rm dual}=\frac{100000}{452349}.
\]
Use the weights in \cref{tab:dual}. There are $53$ positive integers, with
sum $1,000,000$ and largest entry $150,463$. Direct integer addition verifies
all seventeen balance equations \eqref{eq:balance}.

\begin{table}[htbp]
\centering\small
\begin{tabular}{@{}lrl@{}}
\toprule
Row & $r_i$ & Monomial weights, in the order in (\ref{eq:map})\\
\midrule
$c$ & 4,857 & 3,080, 1,777\\
$d$ & 57,691 & 29,588, 28,103\\
$e$ & 242,698 & 93,305, 149,393\\
$f$ & 149,393 & 88,617, 60,776\\
$g$ & 248,519 & 150,463, 98,056\\
$h$ & 47,024 & 27,413, 19,611\\
$p$ & 60,776 & 39,670, 15,071, 3,712, 1,589, 734\\
$q$ & 113,127 & 63,384, 36,448, 7,450, 4,066, 1,779\\
$r$ & 6,370 & 3,666, 2,704\\
$s$ & 27,124 & 18,469, 6,430, 1,247, 679, 299\\
$t$ & 6,546 & 3,753, 2,793\\
$u$ & 11,651 & 7,354, 3,067, 759, 322, 149\\
$v$ & 4,382 & 2,413, 1,560, 289, 120\\
$w$ & 3,026 & 2,033, 796, 139, 58\\
$x$ & 8,357 & 4,786, 2,331, 1,240\\
$y$ & 7,427 & 4,155, 2,506, 766\\
$z$ & 1,032 & 702, 256, 74\\
\bottomrule
\end{tabular}

\caption{An explicit balanced-monomial certificate. The coefficient of every
monomial is one. Within each coordinate row, weights follow the term order
in \eqref{eq:map}.}
\label{tab:dual}
\end{table}

For example, the incoming $C$ exponent is contributed only by the
$\xi c$ terms in rows $Y$ and $Z$. Their weights are $4155$ and $702$,
whose sum is $4857=r_C$. Other balance equations count repeated variables
with multiplicity: the $z^2$ in row $U$ contributes twice its weight to the
incoming $Z$ exponent.

Exact rational interval evaluation proves
\begin{equation}\label{eq:dual-interval}
 \frac{1000668023141660979}{10^{18}}
 \leq \log\mathcal C(\zeta_{\rm dual},m)
 \leq \frac{1000668023155325499}{10^{18}}.
\end{equation}
In particular, $1<\log\mathcal C<1001/1000$, so
$\mathcal C>1$. The slightly generous unit-sized lower margin is merely a
convenient verification statement; scaling all weights scales the logarithm.

\begin{theorem}[Certified limitation of the original system]\label{thm:bracket}
The growth $\mu_{\rm original}$ of the original equality recurrence satisfies
\[
 \frac{452349}{100000}\leq\mu_{\rm original}\leq\frac{9047}{2000}.
\]
Consequently, no supersolution of the original system can establish a growth
upper bound smaller than $4.52349$.
\end{theorem}
\begin{proof}
Balance and \eqref{eq:dual-interval} give the lower bound by
\cref{thm:dual} and \eqref{eq:dual-growth}. The repository's rational
supersolution at $\zeta=2000/9047$, replayed by our verifier without a defect
correction, gives the upper bound. Finally, a supersolution at a smaller
claimed growth value would imply a radius larger than
$100000/452349$, contradicting the certified lower bound.
\end{proof}

The last statement can also be read directly: if a supersolution existed at
some $\zeta'>\zeta_{\rm dual}$, monotonicity of $\Phi$ in its first argument
would make it a supersolution at $\zeta_{\rm dual}$, which the dual
certificate forbids. This is an exact obstruction, not a conclusion drawn
from a numerical optimizer failing to improve the bound.

\subsection{Exact verification without enormous integer products}
Computing the rational product \eqref{eq:dual-product} literally is possible
in principle, but raising many rational factors to large integer weights is
unnecessary. Instead, enclose its logarithm using rational arithmetic.
For $1\leq q\leq2$, let $u=(q-1)/(q+1)$, so $0\leq u\leq1/3$. The identity
\[
 \log q=2\sum_{k\geq0}\frac{u^{2k+1}}{2k+1}
\]
follows by integrating the geometric series for $2/(1-u^2)$ from $0$ to $u$.
The remainder after $M$ terms satisfies
\begin{equation}\label{eq:log-error}
 0\leq \log q-2\sum_{k=0}^{M-1}\frac{u^{2k+1}}{2k+1}
 \leq\frac{2u^{2M+1}}{(2M+1)(1-u^2)}.
\end{equation}
All endpoints are rational when $q$ is rational.

For a general positive rational $q$, write $q=2^e w$ with integer $e$ and
$1\leq w<2$. Evaluate $\log w$ and $\log2$ by
\eqref{eq:log-error}; multiply the interval for $\log2$ by $e$, reversing
endpoint roles when $e<0$. This gives a rational enclosure for $\log q$.

The supplied verifier uses $M=16$. Each individual logarithm interval is
rounded outward to denominator $10^{18}$ before multiplication by its
integer weight. This optional outward rounding prevents huge denominators
when intervals from many factors are added. It cannot invalidate enclosure.
Adding the resulting weighted intervals gives exactly
\eqref{eq:dual-interval}. The same cached enclosure of $\log2$ can be reused;
there are at most $54\times16$ short series terms to evaluate.

The verification cost depends on the bit sizes of the rational monomial
factors, weights, and chosen precision, rather than on the number of iterates
required to witness numerical divergence. No logarithm supplied by a
floating-point library enters the proof. We make no universal claim that
certificate discovery or certificate size is polynomial in the reciprocal
distance to the growth threshold.

\subsection{How the weights were discovered}
The following construction is useful when a finite positive critical point
has been found numerically. Suppose
\[
 \tau=\Phi(\rho,\tau),\qquad
 J=\Phi_{\bX}(\rho,\tau),\qquad
 \ell^TJ=\ell^T,\qquad \ell>0.
\]
Normalize by $C=\sum_i\ell_i\tau_i$ and define
\[
 \beta_{ik}=
 \frac{\ell_i c_{ik}\rho^{b_{ik}}\tau^{\alpha_{ik}}}{C}.
\]
The row sums are $\ell_i\tau_i/C$, and the left-eigenvector identity shows
that the weights are balanced. At $\zeta=\rho$, their logarithmic product
is zero. At $\zeta>\rho$, it becomes
\[
 \left(\sum_{i,k}\beta_{ik}b_{ik}\right)\log(\zeta/\rho)>0
\]
provided at least one positively weighted term contains the size variable.
Rational approximation inside the balance equations then yields integer
weights. This is how the displayed certificate was found.

The approximate critical point and eigenvector are not inputs to the final
verifier. Only the resulting integers and rational evaluation point are.
The completeness proof in \cref{sec:complete} also does not require that a
numerical critical point be correct or even available.

\section{The limits of the exact-prefix hierarchy}
\label{sec:stability}

\subsection{A spectral criterion}
Fix a positive evaluation point $\zeta$ at which the true series $\bA$
converges coordinatewise. Put
\[
 a=\bA(\zeta),\quad p_N=P_N(\zeta),\quad
 J=\Phi_{\bX}(\zeta,a),\quad
 J_N=\Phi_{\bX}(\zeta,p_N),
\]
and define the nonnegative forcing vector
\[
 b_N=\bigl(\tail{N}{\Phi(\xi,P_N(\xi))}\bigr)(\zeta).
\]
At this evaluation point,
\begin{equation}\label{eq:tail-numeric}
 \Psi_N(\zeta,y)=b_N+
       \Phi(\zeta,p_N+y)-\Phi(\zeta,p_N).
\end{equation}
We have $p_N\uparrow a$ and $J_N\to J$. Moreover, $b_N\to0$:
coefficientwise $\Phi(\xi,P_N)\preceq\Phi(\xi,\bA)$, and the latter
series converges at $\zeta$. Its high-degree tail bounds $b_N$.

\begin{theorem}[Strict stability and obstruction]\label{thm:stability}
Assume the preceding setup.
\begin{enumerate}[label=(\alph*)]
\item If $\spr(J)<1$, then all sufficiently large $N$ admit a real upper
certificate at $\zeta$. If the polynomial data and $\zeta$ are rational,
the certificate can be chosen rational.
\item For any $N$, if $J_N$ is irreducible,
$\spr(J_N)\geq1$, and $b_N\ne0$, then no upper certificate at $\zeta$
exists for that $N$.
\item If $\spr(J)>1$, $J_N$ is irreducible for all sufficiently large $N$,
and $b_N\ne0$ for all sufficiently large $N$, then no finite-prefix member
of the hierarchy admits an upper certificate at $\zeta$.
\end{enumerate}
\end{theorem}
\begin{proof}
For (a), the nonnegative matrix geometric series gives
\[
 h=(I-J)^{-1}\bone>0,\qquad Jh=h-\bone.
\]
For sufficiently small $\varepsilon>0$, polynomial expansion in the
nonnegative direction $\varepsilon h$ gives, uniformly for $0\leq p_N\leq a$,
\[
 \Phi(\zeta,p_N+\varepsilon h)-\Phi(\zeta,p_N)
 \leq\varepsilon Jh+C\varepsilon^2\bone
\]
for some finite $C\geq0$. Choose $\varepsilon$ with
$C\varepsilon^2\leq\varepsilon/4$, and then $N$ sufficiently large that
$b_N\leq(\varepsilon/4)\bone$. Equation \eqref{eq:tail-numeric} yields
\[
 \Psi_N(\zeta,\varepsilon h)
 \leq\varepsilon h-(\varepsilon/2)\bone.
\]
Thus $p_N+\varepsilon h$ is a certificate with strict slack and strictly
positive tail budget. When the data and evaluation point are rational,
a sufficiently close rational vector preserves both strict inequalities.

For (b), nonnegative polynomial expansion gives
$\Psi_N(\zeta,y)\geq b_N+J_Ny$ for $y\geq0$.
Let $\ell>0$ be a left Perron vector of the irreducible matrix $J_N$,
with eigenvalue $r=\spr(J_N)\geq1$.
A certificate would give $y\geq b_N+J_Ny$, and hence
\[
 (1-r)\ell^Ty\geq\ell^Tb_N>0,
\]
which is impossible.

For (c), continuity of spectral radius gives $\spr(J_N)>1$ for all
sufficiently large $N$. Part (b) rules out certificates at all those levels.
By \cref{thm:nested}, an earlier certificate would remain a certificate at
all later levels. Therefore no earlier one exists either.
\end{proof}

The equality case $\spr(J)=1$ is deliberately not classified in general.
The nonzero forcing and irreducibility conditions in the obstruction are
not decorative: without them a critical linear map can have supersolutions.
For the Bui map, positive prefix coordinates make the dependency pattern
irreducible, and the $C$ row already supplies a nonzero high-degree forcing
term whenever the top $E$ coefficient is positive.

\subsection{An exact counterexample to asymptotic exactness}
\begin{theorem}[All finite prefixes may retain a growth gap]\label{thm:gap}
There is a scalar proper positive convolution inequality with a true series
of growth $1$ for which the exact-prefix majorants have growth tending to
$3$, not $1$.
\end{theorem}
\begin{proof}
Take
\[
 A(\xi)=\frac{\xi}{1-\xi}=\sum_{n\geq1}\xi^n,
 \qquad \Phi(\xi,x)=\xi+x^2.
\]
The coefficient inequality $A\preceq\xi+A^2$ holds: the coefficient of
$\xi^n$ in $A^2$ is $n-1$ for $n\geq2$.
The exact-prefix defect is
\[
 D_N(\xi)=\sum_{n=3}^N(n-2)\xi^n,
\]
with an empty sum for $N<3$. The corrected equality equation is
\begin{equation}\label{eq:toy}
 B_N=\xi+B_N^2-D_N(\xi).
\end{equation}
Let $c_N(t)=t-D_N(t)$ for real $t$.
Since
\[
 D_\infty(t)=\frac{t^3}{(1-t)^2},
\]
we have $D_\infty(1/3)=1/12$ and
$D_\infty'(t)\leq1$ on $0\leq t\leq1/3$, with equality only at the
right endpoint. A finite truncation has strictly smaller derivative there.
Thus $c_N$ is strictly increasing on this interval,
$c_N(1/4)\leq1/4$, and $c_N(1/3)>1/4$.
There is a unique
\[
 \rho_N\in[1/4,1/3)\quad\hbox{with}\quad c_N(\rho_N)=1/4.
\]
At $t=\rho_N$, the value $v=1/2$ solves \eqref{eq:toy}, and
\[
 P_N(\rho_N)<A(\rho_N)<1/2.
\]
It is therefore a valid upper certificate. For any
$t\in(\rho_N,1/3]$, the discriminant of
$x^2-x+c_N(t)=0$ is negative, so there is no finite real fixed point.
Hence the radius of $B_N$ is exactly $\rho_N$.

The truncated polynomials $D_N$ converge uniformly to $D_\infty$ on
$[0,1/3]$. Every limit point $\rho$ of the sequence $\rho_N$ therefore
satisfies
\[
 \rho-\frac{\rho^3}{(1-\rho)^2}=\frac14,
 \quad\hbox{equivalently}\quad (3\rho-1)^2=0.
\]
Thus $\rho_N\to1/3$, and the corrected growth tends to $3$.
The true series has radius $1$ and growth $1$.
\end{proof}

The same barrier is visible directly in \cref{thm:stability}:
\[
 \Phi_x(t,A(t))=\frac{2t}{1-t}
\]
crosses one at $t=1/3$, well before the true singularity at $t=1$.
Every finite coefficient of $B_N$ eventually agrees with $A$, but radii of
convergence do not pass continuously through this coefficientwise limit.

This counterexample is not a disproof of Bui's separate Question~1 about
increasing neighborhood size with improved geometric recurrences. Here the
local polynomial map is held fixed and only a finite exact prefix is
substituted. Those are different hierarchies. It does show that an eventual
convergence claim for the present prefix method requires a stability argument
or additional geometric refinement.

\subsection{A conditional sensitivity formula for choosing improvements}
A useful guide for future searches is available when a smooth critical branch
has already been justified. Suppose an admissible defect improvement
$D+\varepsilon E$ has critical data $(\rho(\varepsilon),\tau(\varepsilon))$
satisfying
\[
 \Phi(\rho(\varepsilon),\tau(\varepsilon))
 -D(\rho(\varepsilon))-\varepsilon E(\rho(\varepsilon))
 =\tau(\varepsilon).
\]
At $\varepsilon=0$, let $\ell^TJ=\ell^T$ for
$J=\Phi_{\bX}(\rho,\tau)$. Differentiating and multiplying by $\ell^T$
cancels the unknown derivative of $\tau$, yielding
\begin{equation}\label{eq:sensitivity}
 \rho'(0)=
 \frac{\ell^TE(\rho)}
      {\ell^T\bigl(\Phi_\xi(\rho,\tau)-D'(\rho)\bigr)},
\end{equation}
provided the denominator is nonzero.
For a positive tail formulation with $\tau=P(\rho)+y$, the denominator
also equals $\ell^T\Psi_\xi(\rho,y)$, because the difference is
$\ell^T(J-I)P'(\rho)=0$. It is nonnegative, and is positive when some
size-dependent monomial contributes positively.

Formula \eqref{eq:sensitivity} suggests prioritizing verified defect
improvements by their Perron-weighted value. It is conditional on the stated
smooth critical family; it is not used in \cref{thm:main}. In particular,
it never authorizes subtracting an unproved or merely guessed defect.

\section{Verification, formalization, and reproducibility}
\label{sec:verification}

\subsection{The proof dependency chain}
The mathematical chain behind the new upper bound is
\[
\begin{gathered}
 \text{geometric recurrence lemma}
 \ +\ \text{exact marked prefix}\\
 \Longrightarrow\ \text{nonnegative corrected tail system}\\
 \Longrightarrow\ \text{rational supersolution at }500/2249\\
 \Longrightarrow\ \text{coefficient majorant}
 \Longrightarrow\ \lambda\leq2249/500.
\end{gathered}
\]
The new lower certificate for the old equality system follows a different
chain:
\[
\begin{gathered}
 \text{integer balance}\ +\ \text{rational log enclosure}\\
 \Longrightarrow\ \text{no positive supersolution}\\
 \Longrightarrow\ \mu_{\rm original}\geq452349/100000.
\end{gathered}
\]
These chains share the monomial map but not the logical direction of the
majorant comparison.

\subsection{Archive contents and commands}
The archive includes the article source and PDF, the exact marked profile,
twelve upper certificates, the dual certificate, the verification transcript,
and source code. The principal files are:
\begin{center}\small
\begin{tabular}{@{}ll@{}}
\toprule
File & Purpose\\
\midrule
\texttt{code/model.py} & Sparse map, patterns, exact polynomial operations\\
\texttt{code/verify.py} & Standard-library rational certificate replay\\
\texttt{code/enumerate.cpp} & Incremental marked enumeration through size $18$\\
\texttt{code/enumerate\_direct.cpp} & Direct-scanning comparison implementation\\
\texttt{code/research.py} & Independent Python enumeration and numerical exploration\\
\texttt{code/discover\_certificates.py} & Numerical search and rationalization\\
\texttt{data/profiles.json} & Exact marked and unmarked counts\\
\texttt{data/upper\_N18.json} & Main rational upper certificate\\
\texttt{data/dual\_original.json} & Balanced lower certificate for original model\\
\texttt{data/verification.json} & Exact arithmetic residuals and log interval\\
\bottomrule
\end{tabular}
\end{center}

From the archive root, the fast arithmetic audit is
\begin{lstlisting}
python3 code/verify.py --report data/verification.json
\end{lstlisting}
It uses only Python's standard library, arbitrary-precision integers, and
\texttt{fractions.Fraction}. It checks the monomial graph hypotheses,
nonnegative prefix defects, five occupancy partitions, all upper-certificate
residuals and tail budgets, all dual balance equations, the rational log
interval, and the repository's previous $4.5235$ supersolution.

Full regeneration of the main finite table is separate:
\begin{lstlisting}
c++ -O3 -std=c++17 code/enumerate.cpp -o enumerate
./enumerate 18 > regenerated.csv
python3 code/check_profile.py regenerated.csv
\end{lstlisting}
This visits more than two billion partial objects in total and is much slower
than certificate replay. The comparison tool checks all marked and unmarked
entries against the archived table. No external sequence database is needed
for either enumeration or certificate checking.

The numerical discovery scripts require NumPy, SciPy, and SymPy. They are
provided for further experimentation, but are outside the proof boundary.
The figures in the numerical-estimate column of \cref{tab:progress} should
not be promoted to exact assertions merely because those scripts print a
small residual or a successful solver status.

\subsection{A focused Lean extension plan}
The most direct formal extension reuses
\texttt{geometricPublishedBuiRecurrences} and the existing domination of the
fixed-polyomino count by the $G$ coordinate. The genuinely new interfaces are
an exact finite prefix for the marked profile, the defect-corrected
weighted-prefix inequality, and the new rational vector.

A practical route can avoid formalizing complex singularity theory entirely.
First prove the finite convolution identity underlying
\eqref{eq:prefix-def}. Next prove a weighted-tail version of the repository's
monotone-prefix certificate theorem. Instantiate it with the new vector,
using ordinary exact rational reduction for every inequality. The existing
growth endpoint can then carry the pointwise exponential bound to $\lambda$.
The expensive part is a formally verified enumeration or a small proof object
certifying the marked prefix, not the seventeen final inequalities.

For dual certificates, the balance identity is finite integer arithmetic.
One can either formalize weighted arithmetic--geometric mean directly or
expand rational weights into an unweighted finite product argument. A
logarithmic checker with the remainder bound \eqref{eq:log-error} avoids
large powers and is a natural reusable component. The abstract completeness
theorem need not be formalized before the particular $53$-weight certificate
can be checked.

No Lean source claiming these new results is included. In particular, the
presence of a Lean proof of the old bound in the repository must not be
confused with a Lean proof of the new finite-prefix theorem or of the
size-$18$ data.

\section{Proposed research questions and next projects}
\label{sec:questions}

The following projects are intended as concrete continuations, not claims
already established in this report.

\begin{enumerate}[label=\textbf{\arabic*.},leftmargin=2em]
\item \textbf{Kernel-check the $4.498$ bound.}
Formalize the finite-prefix certificate lemma and a verified marked
polyomino enumerator, or design a compact reflective certificate for its
output. The target should be an unconditional endpoint with precisely the
same counting conventions as \texttt{GeometricComplete.lean}. This is the
highest-priority validation project.

\item \textbf{Generate longer marked prefixes without listing every object.}
Adapt a transfer-matrix or connectivity-state enumeration to accumulate all
seventeen occurrence counts, not merely $A_n$. The relevant challenge is
preserving marked local information while merging equivalent boundary
states. Exact counts of larger objects could strengthen the present bound
without changing the imported geometric proof.

\item \textbf{Certify lower bounds for the corrected envelopes.}
Expand or represent the positive maps $\Psi_N$ and construct balanced-monomial
certificates for them. This would replace the middle column of
\cref{tab:progress} by narrow rigorous two-sided brackets. Sparsity or an
arithmetic-circuit representation may avoid a large explicit expansion.

\item \textbf{Locate the actual stability barrier of the fixed Bui map.}
Bound $\spr(\Phi_{\bX}(\zeta,\bA(\zeta)))$ using certified lower and upper
bounds on the true marked generating functions. Determine whether the
prefix hierarchy's limiting obstruction lies substantially below the true
polyomino radius. The strict criteria in \cref{thm:stability} give concrete
sufficient tests, while the equality case remains open.

\item \textbf{Optimize which geometric deficiencies to repair.}
Use Perron-weighted sensitivities as a search heuristic, but accept a repair
only after proving the corresponding injection or coefficient improvement.
A small new neighborhood resolving a heavily weighted overcount may be more
valuable than many additional unweighted states.

\item \textbf{Study Bui's neighborhood-refinement convergence question.}
Develop a hierarchy in which neighborhood size and geometric recurrences
are refined together. Prove convergence to the true growth constant, or
identify a rigorous obstruction. The scalar counterexample here rules out
one naive prefix-only argument; it does not settle the neighborhood question.

\item \textbf{Quantify certificate bit complexity.}
Given a rational target separated from the critical value by a known amount,
bound the sizes of balanced integer weights and the number of logarithm
terms required. Separate complexity of discovery from complexity of replay.
The present million-unit total weight is an example, not a worst-case bound.

\item \textbf{Extend the alternative beyond the irreducible nonlinear case.}
Decompose reducible systems into strongly connected blocks and analyze
linear, unproductive, and boundary blocks separately. A comprehensive theorem
should explain when weak infeasibility can occur and when a strict rational
obstruction still exists. Maximum operators require a different or enriched
certificate language.

\item \textbf{Obtain certified critical asymptotics.}
For selected corrected systems, prove uniqueness of the relevant positive
critical point, nondegeneracy, and the resulting coefficient asymptotics.
Such a result would justify numerical critical continuation as mathematics,
not merely as a source of candidate certificates.

\item \textbf{Transfer the method to other lattice animals and local grammars.}
Identify an existing positive recurrence system on a different lattice,
compute its first exact defects, and test whether a modest prefix correction
produces a useful improvement. The general theorems apply without change
once properness and the underlying combinatorial domination are proved.
\end{enumerate}

\section{Conclusion}
The repository's polyomino project supports more than a small numerical
retuning. Incorporating exact marked prefixes changes the majorant itself,
and the resulting rational certificate proves $\lambda\leq4.498$.
Balanced monomial weights provide the missing lower-certificate mechanism
for a natural class of convolution systems; their explicit use also proves
that the original recurrence cannot support the new bound without a
structural correction. Finally, the spectral analysis identifies why this
successful improvement must not be extrapolated automatically to asymptotic
exactness. The central next step is independent verification, especially
formal certification of the finite marked enumeration.

\appendix
\clearpage
\section{Complete marked profiles through size eighteen}
\label{app:data}

The tables give the exact occurrence counts for the patterns in
\cref{tab:patterns}. All omitted size-zero entries are zero. The row $A_n$
counts translation classes; every other row counts marked occurrences.
These data, together with \eqref{eq:map}, determine every coefficient of
$D_{18}$ by finite integer convolution. The JSON and CSV forms contain the
same integers without typographical formatting.

\begin{table}[H]
\centering\footnotesize
\setlength{\tabcolsep}{4pt}
\renewcommand{\arraystretch}{1.12}
\begin{tabular}{@{}lrrrrrrrrr@{}}
\toprule
$n$ & 1 & 2 & 3 & 4 & 5 & 6 & 7 & 8 & 9\\
\midrule
$A_n$ & 1 & 2 & 6 & 19 & 63 & 216 & 760 & 2725 & 9910\\
\midrule
$C$ & 1 & 1 & 1 & 3 & 10 & 35 & 128 & 479 & 1819\\
$D$ & 1 & 1 & 2 & 6 & 20 & 71 & 260 & 973 & 3696\\
$E$ & 1 & 1 & 3 & 10 & 35 & 128 & 479 & 1821 & 7001\\
$F$ & 1 & 3 & 10 & 35 & 128 & 481 & 1841 & 7131 & 27856\\
$G$ & 1 & 2 & 6 & 20 & 71 & 261 & 982 & 3751 & 14481\\
$H$ & 1 & 2 & 5 & 15 & 51 & 182 & 670 & 2517 & 9592\\
$P$ & 0 & 1 & 4 & 15 & 57 & 220 & 859 & 3380 & 13375\\
$Q$ & 0 & 1 & 3 & 10 & 36 & 133 & 503 & 1930 & 7480\\
$R$ & 0 & 1 & 3 & 8 & 25 & 85 & 303 & 1110 & 4150\\
$S$ & 0 & 1 & 3 & 9 & 31 & 111 & 410 & 1544 & 5896\\
$T$ & 0 & 1 & 3 & 9 & 30 & 106 & 388 & 1451 & 5510\\
$U$ & 0 & 1 & 4 & 14 & 50 & 185 & 700 & 2686 & 10412\\
$V$ & 0 & 0 & 1 & 4 & 14 & 52 & 196 & 748 & 2884\\
$W$ & 0 & 0 & 1 & 4 & 13 & 46 & 168 & 625 & 2362\\
$X$ & 0 & 1 & 2 & 5 & 16 & 54 & 192 & 703 & 2626\\
$Y$ & 0 & 1 & 2 & 4 & 12 & 39 & 135 & 485 & 1788\\
$Z$ & 0 & 1 & 2 & 3 & 8 & 25 & 84 & 296 & 1080\\
\bottomrule
\end{tabular}

\caption{Exact profiles for sizes $1$ through $9$.}
\end{table}

\begin{table}[H]
\centering\footnotesize
\setlength{\tabcolsep}{2.25pt}
\renewcommand{\arraystretch}{1.12}
\begin{tabular}{@{}lrrrrrrrrr@{}}
\toprule
$n$ & 10 & 11 & 12 & 13 & 14 & 15 & 16 & 17 & 18\\
\midrule
$A_n$ & 36446 & 135268 & 505861 & 1903890 & 7204874 & 27394666 & 104592937 & 400795844 & 1540820542\\
\midrule
$C$ & 6979 & 26988 & 105011 & 410629 & 1612207 & 6351215 & 25091789 & 99372469 & 394381157\\
$D$ & 14190 & 54917 & 213865 & 837024 & 3289287 & 12969734 & 51285522 & 203287656 & 807484165\\
$E$ & 27144 & 105927 & 415489 & 1636445 & 6467194 & 25630675 & 101822518 & 405337049 & 1616425239\\
$F$ & 109503 & 432550 & 1715129 & 6821554 & 27199077 & 108673645 & 434959946 & 1743471996 & 6997283445\\
$G$ & 56348 & 220593 & 867716 & 3426361 & 13572506 & 53905206 & 214569185 & 855714279 & 3418228578\\
$H$ & 36929 & 143269 & 559155 & 2192740 & 8632427 & 34094245 & 135023622 & 535974728 & 2131792166\\
$P$ & 53155 & 211957 & 847413 & 3395193 & 13626571 & 54768439 & 220390761 & 887757717 & 3579054867\\
$Q$ & 29204 & 114666 & 452227 & 1789916 & 7105312 & 28274531 & 112746667 & 450377230 & 1801803339\\
$R$ & 15744 & 60366 & 233323 & 907529 & 3547802 & 13926810 & 54857224 & 216706457 & 858185665\\
$S$ & 22739 & 88352 & 345290 & 1355716 & 5343140 & 21124511 & 83738100 & 332687072 & 1324308001\\
$T$ & 21162 & 81966 & 319543 & 1252166 & 4927255 & 19455297 & 77039780 & 305805462 & 1216400095\\
$U$ & 40678 & 159874 & 631249 & 2501630 & 9943629 & 39622046 & 158205688 & 632793092 & 2534808294\\
$V$ & 11212 & 43866 & 172467 & 680772 & 2695907 & 10704713 & 42601646 & 169868638 & 678451341\\
$W$ & 9035 & 34875 & 135550 & 529733 & 2079362 & 8191718 & 32369106 & 128231694 & 509107371\\
$X$ & 9950 & 38100 & 147076 & 571394 & 2231348 & 8750584 & 34438134 & 135936824 & 537948754\\
$Y$ & 6709 & 25491 & 97773 & 377796 & 1468440 & 5735092 & 22488118 & 88474763 & 349078294\\
$Z$ & 4025 & 15214 & 58108 & 223734 & 866978 & 3377002 & 13210216 & 51861498 & 204221945\\
\bottomrule
\end{tabular}

\caption{Exact profiles for sizes $10$ through $18$.}
\end{table}
\clearpage

\section{A minimal exact certificate checker}
\label{app:checker}
The following core is sufficient to check an upper candidate once the map
and true prefix are supplied. The full verifier also checks the input
conventions, all intermediate certificates, the dual weights, and its
rational logarithm enclosure.
\begin{lstlisting}[language=Python]
from fractions import Fraction as Q

# P[i] and D[i] are integer coefficient lists, low degree first.
def evaluate(coefficients, zeta):
    value = Q(0)
    for coefficient in reversed(coefficients):
        value = value * zeta + coefficient
    return value

def check_upper(zeta, values, P, D, polynomial_map):
    prefix = [evaluate(p, zeta) for p in P]
    defect = [evaluate(d, zeta) for d in D]
    image = polynomial_map(zeta, values)
    tail = [v - p for v, p in zip(values, prefix)]
    residual = [v - f + d
                for v, f, d in zip(values, image, defect)]
    if any(t < 0 for t in tail):
        raise ValueError("Candidate lies below the known prefix")
    if any(r < 0 for r in residual):
        raise ValueError("A supersolution inequality failed")
    return tail, residual
\end{lstlisting}

The defect is computed, not guessed. If $P_i(\xi)=\sum_{n=1}^Na_{i,n}\xi^n$,
then for each row $i$ one expands the finitely many monomials in
$\Phi_i(\xi,P)$, discards degrees above $N$, and subtracts $P_i$.
Because all numbers involved are integers before evaluation, this step is
particularly suitable for independent implementations or proof-assistant
reflection.

For the dual checker, \eqref{eq:balance} is a sum of integer vectors.
The $j$th entry of a monomial exponent vector is counted with its full
multiplicity. Omitting the factor two for a square, for example, would
invalidate the certificate. The supplied verifier iterates over repeated
variable indices, so this multiplicity is explicit in its sparse format.

\section{Source audit and claim classification}
\label{app:audit}

\begin{longtable}{@{}p{0.25\textwidth}p{0.67\textwidth}@{}}
\toprule
Claim or ingredient & Status and evidence\\
\midrule
\endhead
Geometric inequalities & Imported from Bui v2, Lemma~4 and Appendix~B; corresponding
repository endpoint source inspected at the pinned commit.\\[0.4em]
Old $4.5235$ upper bound & Existing repository result; rational vector independently
replayed in the new verifier. No new Lean build was performed.\\[0.4em]
Prefix-corrected majorant & Proved in \cref{sec:prefix}, including the nonnegative tail
map and the known-prefix lower-budget condition.\\[0.4em]
New $4.498$ bound & Computer-assisted theorem: exact marked enumeration plus a rational
certificate, with the detailed dependency chain in \cref{sec:verification}.\\[0.4em]
Marked table through $18$ & Generated by exhaustive C++ enumeration. Python agrees through
$11$; direct and incremental C++ scans agree through $16$; standard unmarked
totals agree through $18$. Not kernel-checked.\\[0.4em]
Dual obstruction & Elementary weighted AM--GM proof; particular integer balance and
strict log inequality replayed exactly.\\[0.4em]
Certificate completeness & Proved for productive, strongly connected, nonlinear positive
polynomial systems with rational data. Not claimed for all maximum or signed
recurrences. Based on classical geometric-programming ideas.\\[0.4em]
Numerical critical values & Exploratory estimates only; not certified optimal values and
not used to establish the accepted inequalities.\\[0.4em]
Stability theorem and gap & Written proofs in \cref{sec:stability}. The case of limiting
spectral radius exactly one is not classified generally.\\[0.4em]
Sensitivity formula & Conditional on an already justified differentiable critical
family. Not used for the headline bound.\\[0.4em]
Novelty & The explicit bound and application were not found in the inspected sources.
No claim of an exhaustive priority search or independent refereeing.\\
\bottomrule
\end{longtable}

\clearpage
\begin{thebibliography}{20}
\small
\bibitem{bui}
Vuong Bui, \emph{A convolutional approach to bounding the number of
polyominoes}, arXiv:2511.00461v2, 6 May 2026.
\url{https://arxiv.org/abs/2511.00461v2}.
The relevant locations are Section~3, Question~2, and Section~4,
Lemma~4 with Appendix~B. Version~1 had a different title.

\bibitem{repo}
Vladimir Reshetnikov, \emph{ProveIt}, repository snapshot
\sha.
\url{https://github.com/VladimirReshetnikov/ProveIt}.
Repository sources in the next entries are read at this commit.

\bibitem{reporeadme}
\repo, \emph{Klarner's polyomino growth constant}, project README.
\url{https://github.com/VladimirReshetnikov/ProveIt/blob/e21766d04c2b8a9b2cdba0cd43e563024b3bd1b9/Combinatorics/Polyominoes/KlarnerConstant/README.md}.

\bibitem{repocert}
\repo, \texttt{Support/verify\_certificate.py}, independent rational replay
of the $4.5235$ certificate.
\url{https://github.com/VladimirReshetnikov/ProveIt/blob/e21766d04c2b8a9b2cdba0cd43e563024b3bd1b9/Combinatorics/Polyominoes/KlarnerConstant/Support/verify_certificate.py}.

\bibitem{repopatterns}
\repo, \texttt{Lean/KlarnerConstant/Patterns.lean}, the seventeen required
and forbidden offset sets.
\url{https://github.com/VladimirReshetnikov/ProveIt/blob/e21766d04c2b8a9b2cdba0cd43e563024b3bd1b9/Combinatorics/Polyominoes/KlarnerConstant/Lean/KlarnerConstant/Patterns.lean}.

\bibitem{repocomplete}
\repo, \texttt{Lean/KlarnerConstant/GeometricComplete.lean}, composition of
the geometric recurrences and the existing growth endpoint.
\url{https://github.com/VladimirReshetnikov/ProveIt/blob/e21766d04c2b8a9b2cdba0cd43e563024b3bd1b9/Combinatorics/Polyominoes/KlarnerConstant/Lean/KlarnerConstant/GeometricComplete.lean}.

\bibitem{redelmeier}
D.~Hugh Redelmeier, \emph{Counting polyominoes: Yet another attack},
Discrete Mathematics \textbf{36} (1981), 191--203.
\url{https://doi.org/10.1016/0012-365X(81)90237-5}.

\bibitem{oeis}
The OEIS Foundation, \emph{A001168: Number of fixed polyominoes with $n$
cells}, consulted September 2026.
\url{https://oeis.org/A001168}.
Only the positive-size enumerative values are used for comparison; the
entry's historical prose about growth bounds is not used as current evidence.

\bibitem{duffin}
R.~J. Duffin, E.~L. Peterson, and C. Zener,
\emph{Geometric Programming---Theory and Application}, Wiley, 1967.

\bibitem{boyd}
S. Boyd, S.-J. Kim, L. Vandenberghe, and A. Hassibi,
\emph{A tutorial on geometric programming}, Optimization and Engineering
\textbf{8} (2007), 67--127.
\url{https://web.stanford.edu/~boyd/papers/gp_tutorial.html}.

\bibitem{bv}
S. Boyd and L. Vandenberghe, \emph{Convex Optimization},
Cambridge University Press, 2004.
\url{https://web.stanford.edu/~boyd/cvxbook/}.
\end{thebibliography}
\end{document}
